\documentclass[a4paper,11pt]{article}

% ============================================================
%  Le Collé — Séance 1 : Déterminants, courbes paramétrées
%  © 2026 Ahmed Amine Ouqlani
% ============================================================

% ---------- Polices ----------
\usepackage[utf8]{inputenc}
\usepackage[T1]{fontenc}
\usepackage[sc]{mathpazo}          % texte et mathématiques : Palatino
\usepackage[scaled=0.88]{helvet}   % titres et repères : linéale
\linespread{1.07}
\usepackage{amsmath,amssymb,amsthm,mathrsfs,stmaryrd}
\usepackage{microtype}

% ---------- Mise en page ----------
\usepackage[a4paper,top=2.3cm,bottom=2.3cm,left=2.4cm,right=2.4cm,
            headheight=14pt,headsep=0.8cm,footskip=1.1cm]{geometry}
\setlength{\parindent}{0pt}
\setlength{\parskip}{0.45em plus 0.15em minus 0.1em}
\raggedbottom
\emergencystretch=1.5em
\clubpenalty=10000 \widowpenalty=10000 \displaywidowpenalty=10000 \brokenpenalty=10000

% ---------- Couleurs : un bleu de nuit, un ocre, des gris ----------
\usepackage{xcolor}
\definecolor{Ink}{HTML}{1C2533}
\definecolor{Bleu}{HTML}{1F3A5F}
\definecolor{BleuMoyen}{HTML}{4A6A8F}
\definecolor{BleuClair}{HTML}{9DB2C9}
\definecolor{Fond}{HTML}{F2F4F7}
\definecolor{Ocre}{HTML}{B07A26}
\definecolor{Brique}{HTML}{A23B2A}
\definecolor{Gris}{HTML}{5E6672}
\definecolor{Filet}{HTML}{C5CCD6}

% ---------- Packages ----------
\usepackage{enumitem}
\setlist{itemsep=0.25em,topsep=0.35em,parsep=0pt}
\setlist[itemize]{label=\textcolor{BleuMoyen}{\small\textbullet},leftmargin=1.4em}
\setlist[enumerate]{leftmargin=1.8em,label=\textsf{\small\color{Bleu}\arabic*.}}
\usepackage{booktabs,array,multicol}
\usepackage{needspace,float}
\usepackage[font=small,labelfont={sf,color=Bleu},labelsep=period,skip=6pt]{caption}
\usepackage{tikz}
\usetikzlibrary{arrows.meta,calc,decorations.markings}
\usepackage{pgfplots}
\pgfplotsset{compat=1.17}
\usepackage{tkz-tab}
\usepackage[most]{tcolorbox}
\usepackage{titlesec,titletoc}
\usepackage{fancyhdr,lastpage,etoolbox}
\usepackage[colorlinks=true,linkcolor=Bleu,urlcolor=Bleu,citecolor=Bleu,
            pdftitle={Le Collé — TD Propriétés métriques des courbes (énoncés)},pdfauthor={Ahmed Amine Ouqlani}]{hyperref}

\numberwithin{equation}{section}
\numberwithin{figure}{section}
\numberwithin{table}{section}

% ---------- Titres ----------
% Chapitres (sections) : grand numéro, titre en linéale, filet ocre
\titleformat{\section}[display]
  {\normalfont\sffamily\color{Ink}}
  {{\small\color{Gris}\textls[200]{\MakeUppercase{Chapitre}}}\\[-0.2em]%
   {\fontsize{64}{64}\selectfont\color{BleuClair}\thesection}}
  {0.4em}
  {\fontsize{25}{30}\selectfont\bfseries}
  [\vspace{0.5em}{\color{Ocre}\rule{3.2cm}{1.4pt}}]
\titlespacing*{\section}{0pt}{0pt}{1.6em}
\titleformat{name=\section,numberless}[display]
  {\normalfont\sffamily\color{Ink}}{}{0pt}{\fontsize{25}{30}\selectfont\bfseries}
  [\vspace{0.5em}{\color{Ocre}\rule{3.2cm}{1.4pt}}]
\titlespacing*{name=\section,numberless}{0pt}{0pt}{1.6em}

\titleformat{\subsection}
  {\needspace{5\baselineskip}\normalfont\sffamily\Large\bfseries\color{Ink}}
  {\color{Bleu}\thesubsection}{0.8em}{}
\titlespacing*{\subsection}{0pt}{1.8em}{0.5em}
\titleformat{\subsubsection}
  {\needspace{5\baselineskip}\normalfont\sffamily\large\bfseries\color{Ink}}
  {\color{Bleu}\thesubsubsection}{0.7em}{}
\titlespacing*{\subsubsection}{0pt}{1.6em}{0.4em}
\titleformat{\paragraph}[runin]{\normalfont\sffamily\bfseries\color{Ink}}{}{0pt}{}[.]
\titlespacing*{\paragraph}{0pt}{0.9em}{0.6em}
\setcounter{secnumdepth}{3}
\setcounter{tocdepth}{2}

% ---------- Table des matières ----------
\titlecontents{section}[0pt]{\addvspace{0.7em}\sffamily\bfseries\color{Ink}}
  {\makebox[2.2em][l]{\color{Bleu}\thecontentslabel}}{}{\hfill\contentspage}[\addvspace{0.2em}]
\titlecontents{subsection}[2.2em]{\small\linespread{1}\selectfont}
  {\makebox[2.6em][l]{\color{Gris}\thecontentslabel}}{}{\titlerule*[0.7pc]{.}\contentspage}

% ---------- En-têtes et pieds de page ----------
\pagestyle{fancy}
\fancyhf{}
\renewcommand{\sectionmark}[1]{\markboth{\thesection\quad#1}{}}
\fancyhead[L]{\footnotesize\sffamily\color{Gris}\leftmark}
\fancyhead[R]{\footnotesize\sffamily\scshape\color{Gris}Le Collé}
\fancyfoot[C]{\small\sffamily\color{Ink}\thepage}
\fancyfoot[L]{\scriptsize\sffamily\hyperlink{toc}{\color{Filet}sommaire}}
\renewcommand{\headrulewidth}{0.4pt}
\renewcommand{\headrule}{{\color{Filet}\hrule height\headrulewidth width\headwidth}}
\renewcommand{\footrulewidth}{0pt}
\fancypagestyle{plain}{\fancyhf{}\fancyfoot[C]{\small\sffamily\color{Ink}\thepage}\renewcommand{\headrulewidth}{0pt}}

% ---------- Énoncés ----------
\newtheoremstyle{enonce}{0pt}{0pt}{}{}{\sffamily\bfseries\color{Bleu}}{.}{0.5em}
  {\thmname{#1}\thmnumber{~#2}\thmnote{\ \normalfont\sffamily\color{Ink}(#3)}}
\newtheoremstyle{exemple}{0.9em}{0.9em}{}{}{\sffamily\bfseries\color{BleuMoyen}}{.}{0.5em}
  {\thmname{#1}\thmnumber{~#2}\thmnote{\ \normalfont\sffamily\color{Ink}(#3)}}
\newtheoremstyle{remarque}{0.7em}{0.7em}{}{}{\sffamily\itshape}{.}{0.5em}{\thmname{#1}\thmnote{\ (#3)}}
\theoremstyle{enonce}
\newtheorem{definition}{Définition}[section]
\newtheorem{theoreme}[definition]{Théorème}
\newtheorem{proposition}[definition]{Proposition}
\newtheorem{corollaire}[definition]{Corollaire}
\theoremstyle{exemple}
\newtheorem{exemple}[definition]{Exemple}
\theoremstyle{remarque}
\newtheorem*{remarque}{Remarque}

\tcbset{
  defi/.style={enhanced,breakable,colback=white,colframe=white,boxrule=0pt,arc=0pt,
     borderline west={1.2pt}{0pt}{Bleu},left=11pt,right=4pt,top=3pt,bottom=3pt,
     before skip=1.1em,after skip=1.1em},
  thm/.style={enhanced,breakable,colback=Fond,colframe=Fond,boxrule=0pt,arc=0pt,
     borderline west={2.4pt}{0pt}{Bleu},left=11pt,right=8pt,top=6pt,bottom=6pt,
     before skip=1.2em,after skip=1.2em}}
\tcolorboxenvironment{definition}{defi}
\tcolorboxenvironment{theoreme}{thm}
\tcolorboxenvironment{proposition}{thm}
\tcolorboxenvironment{corollaire}{thm}

\renewcommand{\proofname}{Démonstration}
\renewcommand{\qedsymbol}{\textcolor{BleuClair}{$\blacksquare$}}
\makeatletter
\renewenvironment{proof}[1][\proofname]{\par\pushQED{\qed}\normalfont
  \topsep0.5em\relax\trivlist\item[\hskip\labelsep\sffamily\itshape\color{Gris}#1.]\ignorespaces}
  {\popQED\endtrivlist\@endpefalse}
\makeatother

% ---------- Méthode : filets haut et bas, titre en bandeau ----------
\makeatletter\newcommand{\methlabelfix}{\protected@edef\@currentlabel{\thetcbcounter}}\makeatother
\newtcolorbox[auto counter,number within=section]{methode}[2][]{%
  enhanced,unbreakable,colback=white,colframe=white,boxrule=0pt,arc=0pt,
  borderline north={0.9pt}{0pt}{Bleu},borderline south={0.9pt}{0pt}{Bleu},
  left=4pt,right=4pt,top=8pt,bottom=8pt,before skip=1.4em,after skip=1.4em,
  before upper={\methlabelfix\phantomsection\addcontentsline{met}{metentry}{\protect\makebox[2.6em][l]{\thetcbcounter}#2}%
     {\sffamily\bfseries\small\color{white}\colorbox{Bleu}{\strut\ MÉTHODE~\thetcbcounter\ }}%
     \ \ {\sffamily\bfseries\color{Bleu}#2}\par\smallskip},#1}

% ---------- Repères pédagogiques : un filet de couleur, un mot-clé ----------
\newcommand{\definirbloc}[3]{% nom, couleur, libellé
  \newtcolorbox{#1}[1][]{enhanced,breakable,blanker,
    borderline west={1.4pt}{0pt}{#2},left=10pt,right=2pt,top=1pt,bottom=1pt,
    before skip=1em,after skip=1em,
    before upper={{\sffamily\bfseries\footnotesize\color{#2}\textls[120]{\MakeUppercase{#3}}}\quad},##1}}
\definirbloc{reflexe}{Bleu}{Réflexe}
\definirbloc{astuce}{Ocre}{Astuce}
\definirbloc{attention}{Brique}{Attention}
\definirbloc{enpratique}{BleuMoyen}{En pratique}
\definirbloc{plusloin}{Gris}{Pour aller plus loin}
\definirbloc{strategie}{Bleu}{Stratégie}
\definirbloc{lecture}{BleuMoyen}{Lire le schéma}
% compatibilité avec les parties reprises
\newenvironment{garde}{\begin{attention}}{\end{attention}}
\newenvironment{complementHP}{\begin{plusloin}}{\end{plusloin}}
\newenvironment{interpretation}{\begin{enpratique}}{\end{enpratique}}

% ---------- Exercices, corrigés ----------
\newcounter{exo}[section]
\renewcommand{\theexo}{\thesection.\arabic{exo}}
\newcommand{\niveau}[1]{{\color{Ocre}\footnotesize
  \ifcase#1\or$\bullet\circ\circ\circ\circ$\or$\bullet\bullet\circ\circ\circ$\or$\bullet\bullet\bullet\circ\circ$%
  \or$\bullet\bullet\bullet\bullet\circ$\or$\bullet\bullet\bullet\bullet\bullet$\fi}}
\newcommand{\serie}[1]{\par\needspace{8\baselineskip}\vspace{1.5em}%
  \noindent{\sffamily\bfseries\color{Bleu}#1}\par\nobreak\vspace{-0.4em}%
  \noindent{\color{Filet}\rule{\linewidth}{0.4pt}}\par\nobreak\vspace{0.1em}}
\newenvironment{exo}[3]{%
  \par\needspace{5\baselineskip}\vspace{0.75em}\refstepcounter{exo}\label{e:#1}%
  \def\exolab{#1}%
  \noindent{\sffamily\bfseries\color{Bleu}Exercice~\theexo}\quad\niveau{#2}%
  \hfill{\footnotesize\sffamily\color{Gris}#3}\par\nobreak\vspace{0.15em}}%
  {\par\nopagebreak\noindent\hfill{\scriptsize\sffamily\hyperref[c:\exolab]{\color{Gris}corrigé p.~\pageref{c:\exolab}}}\par}
\newenvironment{corrige}[1]{%
  \par\needspace{6\baselineskip}\vspace{1em}\phantomsection\label{c:#1}%
  \noindent{\sffamily\bfseries\color{Bleu}Exercice~\ref{e:#1}}\quad
  {\scriptsize\sffamily\hyperref[e:#1]{\color{Gris}énoncé p.~\pageref{e:#1}}}\par\nobreak\vspace{0.2em}}%
  {\hfill\textcolor{BleuClair}{$\square$}\par}
\newcommand{\etape}[1]{\par\smallskip\noindent{\sffamily\bfseries\small\color{BleuMoyen}#1.}\ }
\newenvironment{solution}{\par\noindent{\sffamily\bfseries\small\color{BleuMoyen}Solution.}\ }{\hfill\textcolor{BleuClair}{$\square$}\par}
\newcommand{\chercher}{\par\smallskip\noindent{\color{Filet}\rule[0.5ex]{1.2cm}{0.4pt}}\
  {\footnotesize\sffamily\itshape\color{Gris}laisser chercher l'étudiant avant de poursuivre}\par\smallskip}

% ---------- Guide de séance ----------

\newcommand{\marge}[1]{\leavevmode\marginpar{\raggedleft\scriptsize\sffamily\color{Gris}\textls[100]{\MakeUppercase{#1}}}\ignorespaces}
\newcommand{\Etape}[4]{%
  \subsection[Étape #1 --- #3 {\color{Gris}(#2)}]{#3}\label{etape:#1}%
  \noindent{\small\sffamily\color{Gris}{\color{Ocre}\bfseries Étape #1}\quad·\quad #2\quad·\quad #4}\par\medskip}
\newcommand{\question}[1]{\par\smallskip\noindent{\sffamily\bfseries\small\color{BleuMoyen}Question.}\ \textit{#1}\par\smallskip}
\newcommand{\rappel}[1]{\par\smallskip\noindent\phantomsection\addcontentsline{rap}{rapentry}{#1}%
  {\sffamily\bfseries\scriptsize\colorbox{Fond}{\color{Gris}\strut RAPPEL}}\ {\sffamily\small\color{Gris}#1.}\ }
\newcommand{\pourquoi}{\par\noindent{\sffamily\bfseries\small\color{Ocre}Pourquoi ?}\ }

% Ouverture de chapitre : \ouverture{objectifs}{prérequis}{durée}
\newcommand{\ouverture}[3]{%
  \begin{tcolorbox}[enhanced,colback=Fond,colframe=Fond,arc=0pt,boxrule=0pt,left=12pt,right=12pt,top=10pt,bottom=10pt,
     before skip=1.2em,after skip=1.6em]
  \begin{minipage}[t]{0.56\linewidth}
  {\sffamily\bfseries\small\color{Bleu}OBJECTIFS}\par\vspace{0.3em}\small #1
  \end{minipage}\hfill
  \begin{minipage}[t]{0.38\linewidth}
  {\sffamily\bfseries\small\color{Bleu}PRÉREQUIS}\par\vspace{0.3em}\small #2
  \ifx\relax#3\relax\else\par\vspace{0.6em}{\sffamily\bfseries\small\color{Bleu}DURÉE}\par\vspace{0.2em}\small #3\fi
  \end{minipage}
  \end{tcolorbox}}
% Fin de chapitre
\newenvironment{essentiel}{\par\needspace{12\baselineskip}\vspace{1.5em}%
  \begin{tcolorbox}[enhanced,unbreakable,colback=white,colframe=white,boxrule=0pt,arc=0pt,
   borderline north={1.4pt}{0pt}{Ocre},left=2pt,right=2pt,top=8pt,bottom=4pt]
  {\sffamily\bfseries\color{Ocre}L'ESSENTIEL DU CHAPITRE}\par\medskip}{\end{tcolorbox}}

% ---------- Notations ----------
\newcommand{\K}{\mathbb{K}}
\newcommand{\R}{\mathbb{R}}
\newcommand{\C}{\mathbb{C}}
\newcommand{\N}{\mathbb{N}}
\newcommand{\Z}{\mathbb{Z}}
\newcommand{\Mn}[1]{\mathcal{M}_{n}(#1)}
\newcommand{\GLn}[1]{\mathrm{GL}_n(#1)}
\newcommand{\LE}{\mathcal{L}(E)}
\newcommand{\rg}{\operatorname{rg}}
\newcommand{\vect}{\operatorname{Vect}}
\newcommand{\Mat}{\operatorname{Mat}}
\newcommand{\Id}{\operatorname{Id}}
\newcommand{\sh}{\operatorname{sh}}
\newcommand{\ch}{\operatorname{ch}}
\renewcommand{\th}{\operatorname{th}}
\DeclareMathOperator{\argsh}{argsh}
\newcommand{\vv}[1]{\overrightarrow{#1}}
\newcommand{\ve}[1]{\vec{#1}}
\newcommand{\ient}[2]{\llbracket #1,#2\rrbracket}
\newcommand{\dt}[1]{\left|\begin{matrix}#1\end{matrix}\right|}
\newcommand{\dis}{\displaystyle}

% ---------- Figures ----------
\pgfplotsset{trig format plots=rad}
\pgfplotsset{courbe/.style={axis lines=middle,axis equal image,
  axis line style={Gris,-{Stealth[length=5pt]}},
  xlabel={$x$},ylabel={$y$},ticklabel style={font=\scriptsize\sffamily,color=Gris},
  every axis x label/.style={at={(ticklabel* cs:1)},anchor=west,font=\small},
  every axis y label/.style={at={(ticklabel* cs:1)},anchor=south,font=\small},
  samples=300,smooth,width=0.6\linewidth}}
\tikzset{arcetude/.style={Bleu,line width=1.25pt},
         reste/.style={BleuClair,line width=1pt},
         tgte/.style={Ocre,line width=0.9pt},
         asy/.style={Gris,dashed,line width=0.6pt},
         pt/.style={circle,fill=Ink,inner sep=1.2pt}}

% ---------- Navigation ----------
\makeatletter
\newcommand*\l@rapentry{\@dottedtocline{1}{0em}{0em}}
\newcommand*\l@metentry{\@dottedtocline{1}{0em}{0em}}
\newcommand{\listemethodes}{\@starttoc{met}}
\newcommand{\listerappels}{\@starttoc{rap}}
\makeatother
% onglets de chapitre sur le bord droit
\usepackage{eso-pic}
\newif\ifonglets \ongletsfalse
\newcommand{\dessineronglet}{%
  \edef\ongcol{\ifnum\value{section}=3 Ocre\else Bleu\fi}%
  \edef\ongnum{\the\value{section}}%
  \pgfmathsetmacro{\ongy}{3.2+2.3*(\ongnum-1)}%
  \begin{tikzpicture}[remember picture,overlay]
  \fill[\ongcol] ([xshift=-0.75cm,yshift=-\ongy cm]current page.north east)
     rectangle ([yshift=-\ongy cm-1.9cm]current page.north east);
  \node[text=white,font=\sffamily\bfseries\large] at ([xshift=-0.37cm,yshift=-\ongy cm-0.95cm]current page.north east)
     {\ongnum};
  \end{tikzpicture}}
\AddToShipoutPictureBG{\ifonglets\ifnum\value{section}>0 \dessineronglet\fi\fi}
% blocs de navigation et d'enseignement
\definirbloc{commencer}{Gris}{Si vous commencez ici}
\definirbloc{horsprogramme}{Brique}{Hors programme}
\newcommand{\HP}{\,{\sffamily\bfseries\scriptsize\textcolor{Brique}{[HP]}}}
\newcommand{\source}[1]{{\footnotesize\sffamily\color{Gris}\,(#1)}}
\newcommand{\coupdepouce}[1]{\par\vspace{-0.2em}\noindent{\footnotesize\sffamily\color{Gris}%
  \textbf{Coup de pouce} (à donner si l'étudiant bloque) : #1}\par}
\newcommand{\pointarret}[1]{\par\vspace{1em}\noindent{\color{Filet}\rule[0.5ex]{\linewidth}{0.4pt}}\par\nobreak
  \noindent{\footnotesize\sffamily\color{Ocre}\bfseries POINT D'ARRÊT POSSIBLE}\quad{\footnotesize\sffamily\color{Gris}#1}\par\vspace{0.6em}}
\newcommand{\demarche}[1]{\noindent{\sffamily\bfseries\small\color{BleuMoyen}Démarche.}\ \textit{#1}\par\smallskip}

% ---------- Corrigés : idées, astuces, pièges, à retenir ----------
\newcommand{\choix}[1]{\par\noindent{\sffamily\bfseries\small\color{BleuMoyen}Choix de la méthode.}\ #1\par\smallskip}
\newcommand{\gain}[1]{\par\noindent{\sffamily\bfseries\small\color{Ocre}Astuce.}\ #1\par}
\newcommand{\piege}[1]{\par\noindent{\sffamily\bfseries\small\color{Brique}Piège.}\ #1\par}
\newtcolorbox{retenirbox}{enhanced,breakable,colback=Fond,colframe=Fond,arc=0pt,boxrule=0pt,
  borderline west={2pt}{0pt}{Ocre},left=8pt,right=6pt,top=4pt,bottom=4pt,before skip=0.7em,after skip=0.5em}
\newcommand{\retenir}[1]{\begin{retenirbox}{\sffamily\bfseries\footnotesize\color{Ocre}À RETENIR}\quad #1\end{retenirbox}}
\newcommand{\TD}[1]{{\color{Ocre}\bfseries TD\,#1}\ ·\ }
\newcommand{\q}[1]{\par\smallskip\noindent{\sffamily\bfseries\color{Bleu}#1}\ }
\providecommand{\iddots}{\mathinner{\mkern1mu\raise1pt\hbox{.}\mkern2mu\raise4pt\hbox{.}\mkern2mu\raise7pt\hbox{.}\mkern1mu}}

\providecommand{\Ker}{\operatorname{Ker}}\renewcommand{\Im}{\operatorname{Im}}\providecommand{\tr}{\operatorname{tr}}\providecommand{\Sp}{\operatorname{Sp}}

% ================================================================ TD : mise en page propre au livret d'exercices
\fancyhf{}
\fancyhead[L]{\footnotesize\sffamily\color{Gris}TD · Propriétés métriques des courbes · énoncés}
\fancyhead[R]{\footnotesize\sffamily\scshape\color{Gris}Le Collé}
\fancyfoot[C]{\small\sffamily\color{Ink}\thepage\,/\,\pageref{LastPage}}
\fancypagestyle{plain}{\fancyhf{}\fancyfoot[C]{\small\sffamily\color{Ink}\thepage\,/\,\pageref{LastPage}}\renewcommand{\headrulewidth}{0pt}}
\makeatletter
\newcommand{\partie}[3]{% numéro, titre, description
  \par\needspace{10\baselineskip}\vspace{1.4em}\phantomsection\label{partie:#1}\pdfbookmark[0]{Partie #1}{partie.#1}%
  {\sffamily\small\color{Gris}\textls[200]{PARTIE #1}}\par\nobreak\vspace{0.1em}
  {\sffamily\bfseries\LARGE\color{Ink}#2}\par\nobreak\vspace{0.3em}
  {\color{Ocre}\rule{2.6cm}{1.2pt}}\par\nobreak\vspace{0.2em}
  {\small\color{Gris}#3}\par\nobreak\vspace{0.6em}}
\newcommand{\exercice}[7]{% numéro, difficulté, titre, origine, notions, identifiant, page du corrigé
  \par\needspace{8\baselineskip}\vspace{1.1em}\phantomsection
  \def\@currentlabel{#1}\label{e:#6}\pdfbookmark[1]{Exercice #1}{e.#6}%
  \noindent{\sffamily\bfseries\large\color{Bleu}Exercice~#1}\quad\niveau{#2}\hfill{\footnotesize\sffamily\color{Gris}#4}\par\nobreak
  \vspace{0.1em}\noindent{\sffamily\bfseries\color{Ink}#3}\par\nobreak
  \vspace{0.1em}\noindent{\footnotesize\sffamily\color{Gris}Notions : #5.\hfill\mbox{corrigé p.~#7}\par}\nobreak
  \vspace{0.35em}}
\newcommand{\renvoi}[1]{\par\nopagebreak\vspace{0.2em}\noindent\hfill{\scriptsize\sffamily\color{Gris}#1}\par}
\newcommand{\place}{\par\vspace{1.6cm plus 0.6cm minus 0.6cm}}
\newcommand{\correction}[4]{% numéro, titre, identifiant, renvoi
  \par\needspace{7\baselineskip}\vspace{1.3em}\phantomsection
  \def\@currentlabel{#1}\label{e:#3}\pdfbookmark[1]{Exercice #1}{c.#3}%
  {\color{Filet}\hrule height 0.4pt}\vspace{0.5em}
  \noindent{\sffamily\bfseries\large\color{Bleu}Exercice~#1}\quad{\sffamily\bfseries\color{Ink}#2}\hfill{\scriptsize\sffamily\color{Gris}#4}\par\nobreak
  \vspace{0.4em}}
\newcommand{\sommaireligne}[3]{\noindent\hbox to 2.2em{\hfill\sffamily\color{Bleu}#1.}\enspace\begin{minipage}[t]{0.8\linewidth}\raggedright#2\end{minipage}\hfill\pageref{e:#3}\par}
\makeatother
\setlength{\parskip}{0.5em plus 0.15em minus 0.1em}

\begin{document}
\thispagestyle{plain}
{\sffamily\small\color{Gris}\textls[150]{LE COLLÉ · MATHÉMATIQUES PT · CHAPITRE 6}}\par\vspace{0.2em}
{\sffamily\bfseries\fontsize{26}{30}\selectfont\color{Ink}TD --- Propriétés métriques des courbes}\par\vspace{0.3em}
{\sffamily\Large\color{BleuMoyen}Énoncés}\par\vspace{0.4em}
{\color{Ocre}\rule{3.2cm}{1.4pt}}\par\vspace{0.8em}
{\small Ce livret rassemble 24 exercices, classés en quatre parties, du plus direct au plus exigeant.
Il reprend \textbf{tous les exercices du livret} (TD de la classe, exercices d'entraînement, problème type devoir) et les exemples
du cours qui se prêtent à l'entraînement, complétés par des \textbf{exercices complémentaires}, absents du livret.
L'origine de chaque exercice est indiquée en haut à droite : \emph{« Livret, exercice~1.5 · TD~1 de la classe »}, \emph{« Livret, exemple~2.4 »}
ou \emph{« Exercice complémentaire »}. La difficulté va de \niveau{1} à \niveau{5}.
Les corrections sont dans un livret séparé, \texttt{TD\_metrique\_corrections.pdf} ; la page de chaque correction est indiquée en tête de l'énoncé.}

\vspace{0.6em}{\sffamily\bfseries\color{Bleu}Sommaire}\par\vspace{0.2em}{\small\setlength{\parskip}{1.5pt}\setlength{\columnsep}{1.2cm}\begin{multicols}{2}\raggedcolumns
\vspace{0.5em}\noindent\begin{minipage}[t]{0.86\linewidth}\raggedright\leavevmode\sffamily\bfseries\color{Ink}I. Applications directes du cours\end{minipage}\hfill\pageref{partie:I}\par\vspace{0.15em}
\sommaireligne{1}{Courbure d'un cercle, d'une droite}{met-cercle}
\sommaireligne{2}{Longueur de la cardioïde}{ex-cardio}
\sommaireligne{3}{Longueur d'une arche de cycloïde}{m-cycl}
\sommaireligne{4}{Longueur sur la chaînette}{met-chain}
\sommaireligne{5}{Longueur d'un arc de parabole}{met-parab}
\sommaireligne{6}{Courbure d'un graphe}{met-graphe}
\sommaireligne{7}{Cercle de courbure d'un graphe}{met-osc-exp}
\sommaireligne{8}{Une première enveloppe}{met-env-cercle}
\sommaireligne{9}{Une parabole comme enveloppe}{met-env-parab}
\vspace{0.5em}\noindent\begin{minipage}[t]{0.86\linewidth}\raggedright\leavevmode\sffamily\bfseries\color{Ink}II. Exercices classiques et méthodes incontournables\end{minipage}\hfill\pageref{partie:II}\par\vspace{0.15em}
\sommaireligne{10}{Longueur totale de l'astroïde}{m-astro}
\sommaireligne{11}{Courbure de l'ellipse aux sommets}{m-ellipse}
\sommaireligne{12}{Repère de Frenet et développée de la parabole}{ex-devparab}
\sommaireligne{13}{Courbure de la cardioïde}{met-cardio-courb}
\sommaireligne{14}{La spirale logarithmique}{met-spirale}
\sommaireligne{15}{Longueur et courbure d'une boucle}{met-boucle}
\sommaireligne{16}{L'échelle qui glisse}{ex-echelle}
\sommaireligne{17}{Enveloppe d'un segment glissant}{e-segment}
\sommaireligne{18}{L'hyperbole comme enveloppe}{met-env-hyp}
\vspace{0.5em}\noindent\begin{minipage}[t]{0.86\linewidth}\raggedright\leavevmode\sffamily\bfseries\color{Ink}III. Exercices de synthèse et approfondissement\end{minipage}\hfill\pageref{partie:III}\par\vspace{0.15em}
\sommaireligne{19}{Repère de Frenet et développée de la cycloïde}{m-cycl-dev}
\sommaireligne{20}{Les courbes à courbure constante}{met-courbure-cste}
\sommaireligne{21}{Courbure et développée de la chaînette}{met-chain-dev}
\sommaireligne{22}{Développée de l'ellipse comme enveloppe des normales}{e-ellipse}
\sommaireligne{23}{Développée de l'astroïde}{met-astro-dev}
\vspace{0.5em}\noindent\begin{minipage}[t]{0.86\linewidth}\raggedright\leavevmode\sffamily\bfseries\color{Ink}IV. Exercice difficile et type concours\end{minipage}\hfill\pageref{partie:IV}\par\vspace{0.15em}
\sommaireligne{24}{Tractrice et chaînette (type concours)}{met-tractrice}
\end{multicols}}
\clearpage
\partie{I}{Applications directes du cours}{Longueur, abscisse curviligne, courbure d'un cercle et d'un graphe, première enveloppe.}
\exercice{1}{1}{Courbure d'un cercle, d'une droite}{Exercice complémentaire}{abscisse curviligne, repère de Frenet, courbure}{met-cercle}{1}
\begin{enumerate}
\item Pour le cercle $f(t)=(a+r\cos t,\ b+r\sin t)$, $r>0$, déterminer l'abscisse curviligne d'origine $t=0$, le repère de Frenet, la courbure et le centre de courbure.
\item Même question pour une droite $f(t)=A+t\ve u$, $\|\ve u\|=1$.
\end{enumerate}
\place
\exercice{2}{1}{Longueur de la cardioïde}{Livret, exemple 1.2}{longueur, angle moitié}{ex-cardio}{1}
On considère $f(t)=\big((1-\cos t)\cos t,\ (1-\cos t)\sin t\big)$.
\begin{enumerate}
\item Calculer $\|f'(t)\|^2$ et l'écrire à l'aide de $\sin\frac t2$.
\item En déduire la longueur de la cardioïde ($t\in[0,2\pi]$).
\item Que donnerait le même calcul sur $[0,4\pi]$ ? Commenter.
\end{enumerate}
\place
\exercice{3}{2}{Longueur d'une arche de cycloïde}{Livret, exercice 1.1}{longueur, abscisse curviligne, angle moitié}{m-cycl}{2}
Soit la cycloïde $f(t)=(t-\sin t,\ 1-\cos t)$.
\begin{enumerate}
\item Calculer la longueur d'une arche, $t\in[0,2\pi]$.
\item Déterminer l'abscisse curviligne $s$ d'origine $O=M(0)$, orientée dans le sens des $t$ croissants.
\end{enumerate}
\place
\exercice{4}{1}{Longueur sur la chaînette}{Exercice complémentaire}{longueur, abscisse curviligne}{met-chain}{2}
Soit la chaînette $y=\ch x$. Calculer la longueur de l'arc d'extrémités $A(0,1)$ et $M(a,\ch a)$, $a>0$, puis l'abscisse curviligne d'origine $A$.
\place
\exercice{5}{2}{Longueur d'un arc de parabole}{Exercice complémentaire}{longueur, changement de variable}{met-parab}{2}
Calculer la longueur de l'arc de la parabole $y=\frac{x^2}2$ compris entre les points d'abscisses $0$ et $1$.
\place
\exercice{6}{2}{Courbure d'un graphe}{Exercice complémentaire}{courbure, rayon de courbure}{met-graphe}{3}
Soit $g$ une fonction de classe $\mathcal C^2$ sur un intervalle $I$ et $\Gamma$ son graphe, paramétré par $f(x)=(x,g(x))$.
\begin{enumerate}
\item Montrer que la courbure au point d'abscisse $x$ vaut $\gamma(x)=\dfrac{g''(x)}{\big(1+g'(x)^2\big)^{3/2}}$.
\item Application : rayon de courbure de $y=x^2$ en $O$, et de $y=\ln x$ au point d'abscisse $1$.
\end{enumerate}
\place
\exercice{7}{2}{Cercle de courbure d'un graphe}{Exercice complémentaire}{courbure, centre de courbure, cercle de courbure}{met-osc-exp}{3}
On considère le graphe de la fonction exponentielle, au point $A(0,1)$.
\begin{enumerate}
\item Calculer la courbure et le rayon de courbure en $A$ (on pourra utiliser l'exercice~\ref{e:met-graphe}).
\item Déterminer le centre de courbure $I$ et une équation du cercle de courbure en $A$.
\end{enumerate}
\place
\exercice{8}{1}{Une première enveloppe}{Exercice complémentaire}{enveloppe}{met-env-cercle}{3}
Déterminer l'enveloppe de la famille de droites $\mathcal D_t:x\cos t+y\sin t=1$, $t\in\R$. Interpréter.
\place
\exercice{9}{2}{Une parabole comme enveloppe}{Exercice complémentaire}{enveloppe, parabole}{met-env-parab}{3}
Pour $t\in\,]0,1[$, on note $\mathcal D_t$ la droite passant par $P(t,0)$ et $Q(0,1-t)$.
\begin{enumerate}
\item Montrer que $\mathcal D_t$ a pour équation $(1-t)x+ty=t(1-t)$.
\item Déterminer l'enveloppe de la famille $(\mathcal D_t)$, puis montrer qu'elle est contenue dans la courbe d'équation $\sqrt x+\sqrt y=1$.
\item Vérifier que cette enveloppe est contenue dans la parabole $(x-y)^2-2(x+y)+1=0$.
\end{enumerate}
\place
\clearpage
\partie{II}{Exercices classiques et méthodes incontournables}{Repère de Frenet, rayon et centre de courbure, développées, enveloppes de droites.}
\exercice{10}{3}{Longueur totale de l'astroïde}{Livret, exercice 1.2}{longueur, symétries}{m-astro}{4}
Calculer la longueur totale de l'astroïde $f(t)=(\cos^3t,\sin^3t)$, $t\in[0,2\pi]$.
\place
\exercice{11}{3}{Courbure de l'ellipse aux sommets}{Livret, exercice 1.3}{courbure, rayon de courbure}{m-ellipse}{4}
Soit l'ellipse $f(t)=(a\cos t,b\sin t)$, $a>b>0$.
\begin{enumerate}
\item Montrer que $\alpha'(t)=\frac{\det(f',f'')}{\|f'\|^2}$ (où $\alpha$ est l'angle de la tangente), puis calculer la courbure.
\item Calculer le rayon de courbure aux quatre sommets. Où la courbure est-elle maximale ?
\end{enumerate}
\place
\exercice{12}{3}{Repère de Frenet et développée de la parabole}{Livret, exemple 1.7}{repère de Frenet, courbure, centre de courbure, développée}{ex-devparab}{5}
On considère la parabole $f(t)=\left(t,\frac{t^2}2\right)$, $t\in\R$.
\begin{enumerate}
\item Déterminer le repère de Frenet $(M,\ve T,\ve N)$ en $M(t)$, la courbure $\gamma(t)$ et le rayon de courbure $R(t)$.
\item Déterminer le centre de courbure $I(t)$, puis la développée. Quelle particularité présente-t-elle au point qui correspond au sommet de la parabole ?
\end{enumerate}
\place
\exercice{13}{3}{Courbure de la cardioïde}{Exercice complémentaire}{courbure, rayon de courbure}{met-cardio-courb}{5}
On reprend la cardioïde $f(t)=\big((1-\cos t)\cos t,\ (1-\cos t)\sin t\big)$, $t\in\,]0,2\pi[$.
\begin{enumerate}
\item Calculer $\det(f',f'')$.
\item En déduire la courbure et le rayon de courbure en $M(t)$. Où le rayon de courbure est-il maximal ? Que se passe-t-il quand $t\to0$ ?
\end{enumerate}
\place
\exercice{14}{3}{La spirale logarithmique}{Exercice complémentaire}{abscisse curviligne, angle constant, courbure, développée}{met-spirale}{5}
On considère $f(t)=(e^t\cos t,\ e^t\sin t)$, $t\in\R$.
\begin{enumerate}
\item Calculer $\|f'(t)\|$, puis l'abscisse curviligne $s(t)$ d'origine « $t=-\infty$ » (justifier la convergence). Comparer $s(t)$ et la distance $OM(t)$.
\item Montrer que l'angle entre $\vv{OM}$ et la tangente est constant.
\item Calculer la courbure et le centre de courbure. Reconnaître la développée.
\end{enumerate}
\place
\exercice{15}{3}{Longueur et courbure d'une boucle}{Exercice complémentaire}{point double, longueur, courbure}{met-boucle}{6}
On considère la courbe $f(t)=\left(t^2,\ t-\dfrac{t^3}3\right)$, $t\in\R$.
\begin{enumerate}
\item Montrer que la courbe a un unique point double, atteint pour $t=\pm\sqrt3$.
\item Montrer que $\|f'(t)\|=1+t^2$ et calculer la longueur de la boucle.
\item Calculer la courbure en $M(t)$, puis le rayon de courbure en $O$.
\end{enumerate}
\place
\exercice{16}{3}{L'échelle qui glisse}{Livret, exemple 2.2}{enveloppe, astroïde}{ex-echelle}{6}
Une échelle de longueur $L>0$ glisse, son pied sur l'axe $(Ox)$ et son sommet sur l'axe $(Oy)$. On note $\theta\in\left]0,\frac\pi2\right[$ l'angle qu'elle fait avec le sol.
\begin{enumerate}
\item Montrer que la droite qui porte l'échelle a pour équation $x\sin\theta+y\cos\theta=L\sin\theta\cos\theta$.
\item Déterminer l'enveloppe des positions de l'échelle et la reconnaître.
\end{enumerate}
\place
\exercice{17}{3}{Enveloppe d'un segment glissant}{Livret, exercice 2.1}{enveloppe, système $2\times2$}{e-segment}{7}
Pour $t\in\,]0,1[$, on considère la droite $\mathcal D_t$ passant par $A(t,0)$ et $B(0,1-t)$.
\begin{enumerate}
\item Déterminer l'enveloppe de la famille $(\mathcal D_t)$.
\item Montrer qu'elle est contenue dans la courbe $\sqrt x+\sqrt y=1$, et reconnaître cette courbe.
\end{enumerate}
\place
\exercice{18}{2}{L'hyperbole comme enveloppe}{Exercice complémentaire}{enveloppe, tangente}{met-env-hyp}{7}
Pour $t\neq0$, on considère la droite $\mathcal D_t:x+t^2y=2t$.
\begin{enumerate}
\item Déterminer l'enveloppe de la famille $(\mathcal D_t)$.
\item Vérifier que $\mathcal D_t$ est tangente à l'hyperbole $xy=1$. En quel point ?
\item Interprétation : montrer que le point de contact est le milieu du segment de $\mathcal D_t$ découpé par les axes.
\end{enumerate}
\place
\clearpage
\partie{III}{Exercices de synthèse et approfondissement}{Développées calculées de deux façons, courbes classiques.}
\exercice{19}{4}{Repère de Frenet et développée de la cycloïde}{Livret, exercice 1.4}{Frenet, centre de courbure, développée}{m-cycl-dev}{8}
On reprend la cycloïde $f(t)=(t-\sin t,1-\cos t)$, $t\in\,]0,2\pi[$.
\begin{enumerate}
\item Déterminer $\ve T$, $\ve N$, l'angle $\alpha$ et la courbure $\gamma$.
\item Montrer que le centre de courbure est $I(t)=(t+\sin t,\ \cos t-1)$.
\item En déduire que la développée de la cycloïde est une cycloïde, image de la première par une translation.
\end{enumerate}
\place
\exercice{20}{3}{Les courbes à courbure constante}{Exercice complémentaire}{formules de Frenet, centre de courbure}{met-courbure-cste}{9}
Soit $\Gamma$ une courbe de classe $\mathcal C^2$, régulière, paramétrée par son abscisse curviligne $s\mapsto M(s)$, de repère de Frenet $\big(\ve T(s),\ve N(s)\big)$ et de courbure $\gamma(s)$.
\begin{enumerate}
\item On suppose $\gamma$ constante et égale à $\frac1r$, $r>0$. Montrer que le point $I(s)=M(s)+r\ve N(s)$ est fixe. Conclure.
\item On suppose $\gamma$ identiquement nulle. Montrer que $\Gamma$ est contenue dans une droite.
\end{enumerate}
\place
\exercice{21}{3}{Courbure et développée de la chaînette}{Exercice complémentaire}{rayon de courbure, centre de courbure, normale}{met-chain-dev}{9}
Soit la chaînette $f(x)=(x,\ch x)$.
\begin{enumerate}
\item Montrer que le rayon de courbure vaut $R=\ch^2x$.
\item Déterminer le centre de courbure $I(x)$.
\item Soit $N$ le point où la normale en $M$ coupe l'axe $(Ox)$. Montrer que $MI=MN$ et que $I$ est le symétrique de $N$ par rapport à $M$.
\end{enumerate}
\place
\exercice{22}{4}{Développée de l'ellipse comme enveloppe des normales}{Livret, exercice 2.2}{développée, normales, enveloppe}{e-ellipse}{9}
Soit l'ellipse $f(t)=(a\cos t,b\sin t)$, $a>b>0$.
\begin{enumerate}
\item Écrire l'équation de la normale en $M(t)$.
\item En déduire la développée de l'ellipse comme enveloppe des normales.
\item Vérifier la cohérence avec les rayons de courbure aux sommets (exercice~\ref{e:m-ellipse}).
\end{enumerate}
\place
\exercice{23}{4}{Développée de l'astroïde}{Exercice complémentaire}{repère de Frenet, centre de courbure, développée, similitude}{met-astro-dev}{10}
On considère l'astroïde $f(t)=(\cos^3t,\sin^3t)$, $t\in\left]0,\frac\pi2\right[$.
\begin{enumerate}
\item Déterminer le repère de Frenet et le rayon de courbure en $M(t)$.
\item Montrer que le centre de courbure est $I(t)=\big((2-\cos2t)\cos t,\ (2+\cos2t)\sin t\big)$, ou encore $I(t)=\big(\cos^3t+3\cos t\sin^2t,\ \sin^3t+3\sin t\cos^2t\big)$.
\item On pose $u=x+y$ et $v=y-x$. Montrer que $u=(\cos t+\sin t)^3$ et $v=(\sin t-\cos t)^3$. En déduire que la développée est l'image de l'astroïde par une similitude dont on donnera l'angle et le rapport.
\end{enumerate}
\place
\clearpage
\partie{IV}{Exercice difficile et type concours}{Un problème complet reliant longueur, tangente, courbure et développée.}
\exercice{24}{5}{Tractrice et chaînette (type concours)}{Exercice complémentaire}{abscisse curviligne, tangente, courbure, développée}{met-tractrice}{11}
On considère la courbe $\mathcal T$ (tractrice) définie pour $t>0$ par
\[x(t)=t-\th t,\qquad y(t)=\frac1{\ch t}.\]
\begin{enumerate}
\item Calculer $f'(t)$ et montrer que $\|f'(t)\|=\th t$. En déduire l'abscisse curviligne $s(t)$ d'origine $t=0$ ($M(0)=(0,1)$).
\item Montrer que la tangente en $M(t)$ coupe l'axe $(Ox)$ en $P(t,0)$ et que $MP=1$. Interpréter mécaniquement (« courbe du chien »).
\item Déterminer le repère de Frenet, puis montrer que le centre de courbure est $I(t)=(t,\ch t)$. Reconnaître la développée de la tractrice.
\item Exprimer $y$ en fonction de $s$ et retrouver le rayon de courbure à partir de la relation $R=\frac{ds}{d\alpha}$.
\end{enumerate}
\place
\end{document}
